\documentclass[a4paper]{scrartcl}
\usepackage[utf8]{inputenc}
\usepackage{amsmath,amssymb,amsthm}
\usepackage{enumitem,setspace}
\usepackage[usenames,dvipsnames,table]{xcolor}
\usepackage{graphicx}
\DeclareGraphicsExtensions{.pdf,.png,.eps,.jpg}
\usepackage{algorithm2e}

% --- math operators ---
\usepackage{mathtools}
\DeclarePairedDelimiter\set{\{}{\}} % use $\set{1, 2, 3}$
\DeclarePairedDelimiter\abs{\lvert}{\rvert}
\DeclarePairedDelimiter\norm{\lVert}{\rVert}
\DeclarePairedDelimiter\ceils{\lceil}{\rceil}
\DeclarePairedDelimiter\floor{\lfloor}{\rfloor}
\def\then{\ensuremath{\Rightarrow}} % =>
\def\iff{\ensuremath{\Leftrightarrow}} % if and only if <=>
\def\to{\ensuremath{\rightarrow}} % ->
\def\ot{\ensuremath{\leftarrow}} % ->
\def\Oh{\ensuremath{\mathcal{O}}} % big O like $\Oh(n)$
% ---
\DeclareMathOperator{\preorder}{preorder}
\DeclareMathOperator{\postorder}{postorder}
\DeclareMathOperator{\depth}{depth}

% --- FILL OUT THESE 3 FIELDS ---
\def\sheetNumber{11}
\def\names{Matthias Janßen, Raphael Thelen, Tillmann
Faust} 
\def\sumPoints{20} 

% --- ### BEGIN DOCUMENT ### ---
\begin{document} 

\begin{doublespace} % header
\noindent \textbf{\LARGE{Übungsblatt \sheetNumber}}
\hfill  {\large Punkte: $\boxed{\qquad  /\; \sumPoints}$}\\
{\Large \textbf{Fortgeschrittene Algorithmen (Wintersemester 2022/23)}\\
Bearbeitet von: \textbf{\names}\\
\today}
\end{doublespace}


\section*{Aufgabe 1}
\subsection*{a)}
Seien aktuell \texttt{$S_1.size()+S_2.size()=n$}. Wenn $S_2$ mindestens ein Element enthält dauert die Operation $\Theta(1)$. Wenn $S_2$ leer ist müssen zuerst alle Elemente auf $S_2$ verschoben werden und anschließend das oberste entfernt werden. Da die Operation im worst case maximal $n$ werte bewegt und jede Push- und Pop-Operation für eine Bewegung in $\Theta(1)$ läuft ist die worst-case Laufzeit $\Theta(n)$.

\subsection*{b)}
Per Buchhaltermethode: Wir geben bei jeder Enqueue-Operation 2 Konsteneinheiten zusätzlich zu den Push-Kosten hinzu. Die Gesamtkosten sind damit $\Theta(1)$ und jedes Element hat 2 Kosten als Budget. Wenn nun Dequeue aufgerufen wird liegt entweder ein Element auf $S_2$ und wie haben Kosten 1, oder wir können die Kosten für die Push- und Pop-Operation jeden Elements auf $S_1$ durch die 2 zurückgelegten Kosten bezahlen und müssen nur noch 1 für die Letzte Pop-Operation Zahlen. Damit kostet jede Dequeue-Operation amortisiert nur noch $\Theta(1)$.

\section*{Aufgabe 2}
\subsection*{a)}
\begin{verbatim}
inc() {
	for all i in 0...k-1 {
		if A[i] == 0 {
			A[i] = 1
			return
		} else {
			if i == k-1 {
				set all to 0
				return
			}
			A[i] = 0
		}
	}
	return A
}
\end{verbatim}

\subsection*{b)}
Im worst case sind alle Bits 1 und es muss jedes Bit auf 0 gesetzt werden. Somit sind die worst-case Kosten $\Oh(k)$.

\subsection*{c)}
Wir zahlen für das setzen eines Bits von 0 auf 1 2 €. 1 € wird für den eigentlichen Setzvorgang verwendet, der andere € wird als Kredit gespeichert. Die tatsächlichen Kosten für das Setzen eines Bits auf 1 sind $\Oh(1)$, die amortisierten Kosten liegen bei $\Oh(1)+1$. 
    Somit bleibt jedem Bit 1 Euro Kredit welcher verwendet werden kann, um das Bit später wieder auf 0 zu setzen. Jedes Bit zahlt die Operation selbst. Der Kontostand wird also nie negativ. Somit sind die amortisierten Kosten pro Inkrementierung $\Theta(1)$.

\section*{Aufgabe 1}
\subsection*{a)}
Seien aktuell \texttt{$S_1.size()+S_2.size()=n$}. Wenn $S_2$ mindestens ein Element enthält dauert die Operation $\Theta(1)$. Wenn $S_2$ leer ist müssen zuerst alle Elemente auf $S_2$ verschoben werden und anschließend das oberste entfernt werden. Da die Operation im worst case maximal $n$ werte bewegt und jede Push- und Pop-Operation für eine Bewegung in $\Theta(1)$ läuft ist die worst-case Laufzeit $\Theta(n)$.

\subsection*{b)}
Per Buchhaltermethode: Wir geben bei jeder Enqueue-Operation 2 Konsteneinheiten zusätzlich zu den Push-Kosten hinzu. Die Gesamtkosten sind damit $\Theta(1)$ und jedes Element hat 2 Kosten als Budget. Wenn nun Dequeue aufgerufen wird liegt entweder ein Element auf $S_2$ und wie haben Kosten 1, oder wir können die Kosten für die Push- und Pop-Operation jeden Elements auf $S_1$ durch die 2 zurückgelegten Kosten bezahlen und müssen nur noch 1 für die Letzte Pop-Operation Zahlen. Damit kostet jede Dequeue-Operation amortisiert nur noch $\Theta(1)$.

\section*{Aufgabe 2}
    a)
    \begin{verbatim}
        inc() {
            for all i in 0...k-1 do
                if A[i] == 0
                    A[i] = 1
                    return
                else
                    if i == k-1
                        set all to 0
                        return
                    A[i] = 0
            return A
        }
    \end{verbatim}

    b) Im worst case sind alle Bits 1 und es muss jedes Bit auf 0 gesetzt werden. Somit sind die worst-case Kosten $\Oh(k)$.
\\\\
    c) Wir zahlen für das setzen eines Bits von 0 auf 1 2 €. 1 € wird für den eigentlichen Setzvorgang verwendet, der andere € wird als Kredit gespeichert. Die tatsächlichen Kosten für das Setzen eines Bits auf 1 sind $\Oh(1)$, die amortisierten Kosten liegen bei $\Oh(1)+1$. 
    Somit bleibt jedem Bit 1 Euro Kredit welcher verwendet werden kann, um das Bit später wieder auf 0 zu setzen. Jedes Bit zahlt die Operation selbst. Der Kontostand wird also nie negativ. Somit sind die amortisierten Kosten pro Inkrementierung $\Theta(1)$.

\section*{Aufgabe 3}

\subsection*{a)}

\begin{enumerate}
    \item \textbf{Konstruktion:} Nacheinander $n$ INSERT-Operationen in einen leeren Fibonacci-Heap. $\Rightarrow$ Wurzelliste besteht aus $n$ Bäumen mit Grad 0. 
    \item \textbf{EXTRACTMIN():}
    \begin{itemize}
        \item Minimum entfernen, es verbleiben $n-1$ Knoten.
        \item CLEANUP-Operation.
        \item Konsolidierung durchläuft die gesamte Wurzelliste, echten Kosten sind proportional zur Anzahl der Wurzeln + maximaler Grad, hier jeweils Null.
        \item Mit $n-1$ Wurzeln betragen die echten Kosten
        $$T = O(n - 1 + D(n)) = O(n - 1 + 0) = \Omega(n)$$.
    \end{itemize}
\end{enumerate}

\subsection*{b)}

\begin{enumerate}
    \item \textbf{Konstruktion:} Baum mit Wurzel $w$ und Grad $k$.
    \item Die Wurzel $w$ besitzt $k$ Kinder $c_1, c_2, \dots, c_k$. Wir entfernen für jedes Kind $c_i$ alle Kinder mittels DECREASE-KEY und anschließendem CUT.
    \item Da wir bei DECREASE-KEY nicht markieren, wird $c_i$ nicht wie üblich von $w$ abgeschnitten sondern bleibt teil des Baums.
    \item Sobald alle Enkel entfernt wurden, sind $c_1, \dots, c_k$ Blätter mit Grad 0.
    \item Wurzel $w$ hat weiterhin grad $k$, die Gesamtzahl an Knoten beträgt
    $$1 + k = \Theta(k)$$
\end{enumerate}
\section*{Aufgabe 4}


% -- code for figures
% \begin{figure}[htb]
% \centering
% \includegraphics{filename} % use option [width=\linewidth] to change size to linewidth (or 0.6\linewidth for example)
% \label{fig:name}
% \end{figure}


\end{document}
% --- ### END DOCUMENT ### ---